% I skabelonen er der tre primære ting at ændre: linjerne øverst, brødteksten og teksten i højre kolonne.

% VIGTIGT! For at dokumentet kan compile, skal man benytte XeLaTeX eller LuaLaTeX som compiler. Dette kan gøres i Overleaf ved at gå til Menu -> Settings -> Compiler -> Vælg XeLaTeX/LuaLaTeX

% Brugerspecifikke variable fra YAML-header 
\newcommand{\Titel}{$titel$}   % Vedrørende
\newcommand{\Sagsbehandler}{$sagsbehandler$}   % Sagsbehandler
\newcommand{\Email}{$email$} % Email
\newcommand{\Koncern}{$koncern$}     % Koncern
\newcommand{\Enhed}{$enhed$} % Enhed
\newcommand{\Reference}{$reference$}   % Reference-ID XX ???

\documentclass[12pt]{article}

\usepackage{geometry}
\usepackage{xcolor}
\usepackage{microtype}
\usepackage[skins,breakable]{tcolorbox}
\usepackage{paracol}
\usepackage{setspace}
\usepackage{fontspec}
\usepackage{hyperref}
\hypersetup{
    colorlinks=true,
    linkcolor=blue,
    urlcolor=blue,
    }
\usepackage{array}
\usepackage{booktabs}
\usepackage{footnote}
\usepackage{longtable} % Nødvendig for at kunne inkludere tabeller 
\makeatletter
\def\LT@err#1#2{}
\makeatother
\usepackage{caption}
\providecommand{\tightlist}{%
  \setlength{\itemsep}{0pt}%
  \setlength{\parskip}{0pt}%
}
\providecommand{\pandocbounded}[1]{#1}
$header-includes$
\usepackage{caption}
\usepackage{subcaption}
\usepackage{enumitem}
\setlist[itemize]{noitemsep, topsep=0pt, partopsep=0pt, parsep=1pt, left=0pt}
\setcounter{secnumdepth}{0}
\setlist[itemize,2]{label=\textopenbullet}

\usepackage{polyglossia}
\setdefaultlanguage{danish} % Dato på dansk
\usepackage{ragged2e}

%\makesavenoteenv{longtable}

\newgeometry{top=0.7in, bottom=0.7in, left=0.8in, right=0.8in}


\definecolor{KUrod}{RGB}{144,26,30}

\setmainfont{Times New Roman}

\begin{document}

\begin{flushleft}
    \textls[300]{\vspace*{-0.7cm}\\
    \fontsize{11}{11}\selectfont KØBENHAVNS UNIVERSITET}
\end{flushleft}

\setlength{\columnsep}{0.5in}
\setcolumnwidth{4in, 1.07in}

\begin{paracol}{2}
    \begin{tcolorbox}[blanker, breakable, width=\linewidth]
    \vspace*{0.25in}
    \onehalfspacing
    % Her ændres teksten til venstre for lopoet
    \end{tcolorbox}
    \switchcolumn
    \begin{tcolorbox}[blanker, breakable, width=\linewidth]
        \hfill\includegraphics[width=1.27in]{S:/FAD-REKSEK-REKSTAB-ABI-Analyse/quarto/ku_skabeloner/KU-logo.png}
    \end{tcolorbox}
\end{paracol}

\vspace*{-0.93cm}
\noindent\textcolor{KUrod}{\makebox[\linewidth]{\rule{\paperwidth}{1.05pt}}}   
\vspace*{1.05cm}

\begin{paracol}{2}
    \begin{tcolorbox}[blanker, breakable, width=\linewidth]

    % Her ændres brødteksten (teksten i venstre kolonne)
    \textls[300]{\vspace*{-0.7cm}\\
    \fontsize{10}{11}\selectfont \fontspec{Arial} \textbf{SAGSNOTAT}} \\    
    
    \vspace{0.7 cm}
    \begin{tabular}{@{}p{3cm} p{10cm}@{}}
    \textbf{Vedr.} & \Titel \\
    \textbf{Sagsbehandler} & \Sagsbehandler
    \end{tabular}
    \vspace{0.5 cm}
    \RaggedRight

$body$
        
    \end{tcolorbox}
    \switchcolumn
    \begin{tcolorbox}[blanker, breakable, width=\linewidth+2cm]
    {\fontsize{7}{14}\selectfont
    \fontspec{Arial}

    % Her ændres teksten i højre kolonne
    \today
    \\ \\
    \Koncern \\
    \Enhed \\
    \\ 
    KRYSTALGADE 25\\
    1172 KØBENHAVN K \\
    \\ 
    \\ 
    \Email \\
    \\ \\
    %REF: \Reference

    }
    \end{tcolorbox}
\end{paracol}

\end{document}
